\subsection{First-order Dirac block and protected Pfaffian}
\label{subsec:dirac-pfaffian}

The fourth entry of the Dirac--Igusa pro-object is the protected
Pfaffian line and its first-order Dirac block decomposition.  The
Pfaffian, not the determinant, is the object of interest for a precise
reason: the automorphic target $\Delta_5$ is the square root of the
scalar weight-ten form $\Delta_5^2=\chi_{10}$, and the square root
carries the order-two Maass multiplier $\nu_{\Delta_5}$ that the square
sends to the trivial character.  Geometrically the determinant line
$\mathscr L_{\det}$ is canonically trivialisable and forgets
orientation; its square root, the Pfaffian line, retains the
orientation character $\epsilon_o$, and it is precisely the match
$\epsilon_o=\nu_{\Delta_5}$ that distinguishes the first-order object
$\Delta_5$ from its orientation-blind square $\Delta_5^2$.  Two
definitions of this Pfaffian meet.  The geometric definition starts from
the cosection-reduced d-critical heart, carries a Joyce--Upmeier
orientation, descends to the reduced quotient, and associates a
Pfaffian line to the determinant of the reduced obstruction theory in
the Brav--Bussi--Dupont--Joyce--Szendr\H{o}i framework
\cite{BBDJS2015}.  The algebraic construction starts from the imported
Borcherds--Kac--Moody denominator algebra, fixes a positive window in
$\Gamma_{\mathrm{gram}}$, and writes the algebraic super-Pfaffian
\[
\operatorname{sPf}^{\mathrm{alg}}_R(D)
=\prod_{\gamma\in A_R}(1-x^\gamma)^{\sdim V_\gamma}
\]
of the trivial first-order block decomposition $D_\gamma$ on the
target side.  The Pfaffian--Dirac theorem of Chapter~6 is the
identification of these two Pfaffian sections after the chosen
line-comparison $\iota_{\mathrm{aut}}$ and the four-obstruction
datum of Appendix~\ref{app:obstruction-discharges}.

\subsubsection*{Pfaffian line and squaring isomorphism}

\begin{definition}[Pfaffian line on the cosection-reduced quotient]
\label{def:pfaffian-line}
Let $K^{\mathrm{red}}_{\mathcal C}$ denote the cosection-reduced
self-Ext complex on a charge stratum, with reduced determinant
$\det K^{\mathrm{red}}_{\mathcal C}$ and Joyce--Upmeier orientation
$o_{\mathcal C}$.  The Pfaffian line at finite stage is
\[
\mathscr L_{\mathrm{Pf},R}=
\bigotimes_{\mathcal C}\operatorname{Pf}(K^{\mathrm{red}}_{\mathcal C,R})
\]
on the reduced moduli, with squaring isomorphism
\[
\mathscr L_{\mathrm{Pf},R}^{\otimes 2}\simeq\mathscr L_{\det,R}
\]
in the Brav--Bussi--Dupont--Joyce--Szendr\H{o}i sense
\cite{BBDJS2015}.  The successor maps
$\rho^{\mathrm{Pf}}_{R^+R}$ commute with these squaring isomorphisms,
identifying the inverse limit with a Pfaffian line
$\mathscr L_{\mathrm{Pf},X}$ on the inverse-limit object.  A
normalized section $\operatorname{pf}_{X,R}$ at finite stage and
$\operatorname{pf}_X=\varprojlim_R\operatorname{pf}_{X,R}$ are part
of the datum.
\end{definition}

\subsubsection*{Pfaffian-to-automorphic comparison}

The recognition statement is read on the automorphic line.  A
section-level identification $\operatorname{pf}_X=\Delta_5$ requires
a comparison of $\mathscr L_{\mathrm{Pf},X}$ with the pullback of the
automorphic weight-five line to the moduli base.

\begin{definition}[Pfaffian-to-automorphic line comparison]
\label{def:pfaffian-to-automorphic-line-comparison}
A section-level Pfaffian-to-automorphic comparison is a chosen
isomorphism of line bundles
\[
\iota_{\mathrm{aut}}:\mathscr L_{\mathrm{Pf},X}
\xrightarrow{\sim}\mathcal L^5\otimes\nu_{\Delta_5}
\]
to the pullback of the automorphic weight-five line, twisted by the
multiplier $\nu_{\Delta_5}$, compatible with cusp trivialization at
$\mathfrak c_\infty$ and with the squaring map
$\mathscr L_{\mathrm{Pf},X}^{\otimes 2}\simeq\mathscr L_{\det,X}$.
\end{definition}

\begin{definition}[Finite geometric Pfaffian datum
\((P_{\mathrm{fin}})\)]
\label{def:finite-geometric-pfaffian-datum}
At finite HN height \(R\), a finite geometric Pfaffian datum is the
following system on the reduced compact quotient
\(\mathfrak Q_R=(\mathfrak M_R/E)^{\mathrm{red}}\).
\begin{enumerate}
\item[\textup{(P1)}] For every retained stratum \(\mathcal C\), a
cosection-reduced perfect self-Ext complex
\(K^{\mathrm{red}}_{\mathcal C,R}\) with odd skew self-duality
\[
\vartheta_{\mathcal C,R}:K^{\mathrm{red}}_{\mathcal C,R}
\xrightarrow{\sim}
(K^{\mathrm{red}}_{\mathcal C,R})^\vee[1],
\]
compatible with the reduced obstruction theory and the quotient by
\(E\).
\item[\textup{(P2)}] A Joyce--Upmeier orientation
\(o_{\mathcal C,R}\) and the corresponding Pfaffian line
\[
\mathscr L_{\mathrm{Pf},R}
=\bigotimes_{\mathcal C}\operatorname{Pf}
(K^{\mathrm{red}}_{\mathcal C,R}),
\qquad
\mathscr L_{\mathrm{Pf},R}^{\otimes 2}\simeq\mathscr L_{\det,R}.
\]
\item[\textup{(P3)}] A finite-rank skew first-order block
\(D_{R,\gamma}\) on each retained primitive degree
\(\gamma\in\Gamma_R^{\Pi,+}\), with Pfaffian section
\[
\operatorname{pf}_{X,R}\in
H^0(\mathfrak Q_R,\mathscr L_{\mathrm{Pf},R}).
\]
In a cusp frame this section has support exactly on the retained
active degrees, full even/odd parity dimensions
\(p^{\bar0}_{R,\gamma},p^{\bar1}_{R,\gamma}\), and signed exponent
\(p^{\bar0}_{R,\gamma}-p^{\bar1}_{R,\gamma}=a_\Delta(\gamma)\).
\item[\textup{(P4)}] For every retained simple type-II wall chart,
with wall coordinate \(u_{\delta_i,R}\), a local factorization
\[
\operatorname{pf}_{X,R}
=\upsilon_{\delta_i,R}\,u_{\delta_i,R}\cdot
\operatorname{pf}^{\mathrm{tan}}_{i,R},
\qquad
\operatorname{ord}_{u_{\delta_i,R}=0}(\operatorname{pf}_{X,R})=1,
\]
where the unit \(\upsilon_{\delta_i,R}\) is invariant under the chosen
reflection lift and the normal Pfaffian rank is
\(N^{\mathrm{Pf}}_{\delta_i,R}=1\).  No type-I wall is allowed in this
datum.
\item[\textup{(P5)}] A finite-stage comparison
\(\iota_{\mathrm{aut},R}:\mathscr L_{\mathrm{Pf},R}\to
\mathcal L^5\otimes\nu_{\Delta_5}\) on the automorphic chart,
compatible with the cusp trivialization, the leading coefficient
\(\operatorname{lc}_{\mathfrak c_\infty}(\operatorname{pf}_{X,R})=64\),
and squaring to the determinant line.
\item[\textup{(P6)}] Strict transition isomorphisms
\(\rho^{\mathrm{Pf}}_{R^+R}:
\mathscr L_{\mathrm{Pf},R^+}|_{\mathfrak Q_R}
\xrightarrow{\sim}\mathscr L_{\mathrm{Pf},R}\) carrying
\(\operatorname{pf}_{X,R^+}\) to \(\operatorname{pf}_{X,R}\), with
closed images and vanishing \(\lim^1\) for the line and section
systems.
\end{enumerate}
The scalar Borcherds product, the squared determinant
\(\Delta_5^2\), and the OP scalar \(-4096\Delta_5^{-2}\) are not
entries of \((P_{\mathrm{fin}})\).  They are shadows obtained after
choosing the cusp frame and, in the OP case, after squaring and
inverting the Pfaffian section.

After finite homogeneous trivialisations, \((P_{\mathrm{fin}})\) is
recorded by the Pfaffian proof tables
\[
\mathrm{strata},\quad
\mathrm{skew\_complexes},\quad
\mathrm{pfaffian\_lines},\quad
\mathrm{pfaffian\_sections},\quad
\mathrm{wall\_charts},\quad
\mathrm{automorphic\_comparisons},\quad
\mathrm{transitions},\quad
\mathrm{scalar\_firewall}.
\]
The rows are not scalar product rows.  They record retained strata,
cosection-reduced skew perfect complexes, Pfaffian lines and their
squares, finite Pfaffian sections, type-II rank-one wall charts,
Pfaffian-to-automorphic line comparisons, strict transition maps, and
the exclusion of scalar Borcherds products, squared determinants, OP
scalar traces, and exponent tables as entries of \((P_{\mathrm{fin}})\).
Defect ranks, exponent residuals, leading residuals, and
\(R^1\varprojlim\)-ranks must be zero; the type-II normal Pfaffian rank
and divisor order must be one; the simple-reflection sign is \(-1\);
the automorphic comparison has leading coefficient \(64\), weight \(5\),
and character \(\nu_{\Delta_5}\).  Every row must carry geometric
source provenance and a proof reference.
\end{definition}

\begin{proposition}[Scalar Pfaffian data do not determine a Pfaffian section;
\textit{proved}]
\label{prop:scalar-pfaffian-data-not-section}
Let \(U\) be a quotient chart in the type-II wall complement, let
\(\mathcal L^5\) be the automorphic weight-five line on \(U\), and let
\(T_\chi\) be a non-trivial order-two flat line with character
\(\chi:\pi_1(U)\to\{\pm1\}\).  A divisor, a leading coefficient, a
scalar Borcherds product written after choosing a cusp trivialization of
\(\mathcal L^5\), or the square of that product in \(\mathcal L^{10}\)
does not determine a section of \(\mathcal L^5\otimes T_\chi\).  Thus a
section-level Pfaffian identity
\[
\operatorname{pf}_X=\Delta_5
\]
requires, in addition to the scalar product data, the Pfaffian line
\(\mathscr L_{\mathrm{Pf},X}\), the section \(\operatorname{pf}_X\), the
comparison \(\iota_{\mathrm{aut}}\), the order-two character
\(\epsilon_o=\nu_{\Delta_5}\), and the inverse-limit transition data for
these objects.  In particular, the exponent table \(f(nm,l)\), the
squared identity \(\Delta_5^2=\chi_{10}\), and the OP scalar
\(-4096\,\Delta_5^{-2}\) do not construct
\(\operatorname{pf}_{X,R}\) or the orientation character
\(\epsilon_{o,R}\).
\end{proposition}

\begin{proof}
Choose a simply connected open set \(V\subset U\) and a flat frame
\(\tau\) of \(T_\chi|_V\).  For a local section \(s\) of \(\mathcal L^5\),
the section \(s\otimes\tau\) of \(\mathcal L^5\otimes T_\chi\) has the
same local scalar expression as \(s\) after the frame \(\tau\) is set
equal to \(1\).  Under the canonical trivialization
\(T_\chi^{\otimes2}\simeq\mathcal O_U\), the squares of \(s\) and
\(s\otimes\tau\) give the same section \(s^2\) of \(\mathcal L^{10}\).
Their divisors and leading coefficients in the chosen cusp
trivialization are therefore identical.

The two objects are nevertheless different globally.  Along a loop
\(\ell\) with \(\chi(\ell)=-1\), analytic continuation fixes \(s\) and
sends \(s\otimes\tau\) to \(-s\otimes\tau\).  This monodromy is exactly
the information carried by the order-two orientation character.  Hence
the square, the local product, and the scalar trace forget the data
needed to define the Pfaffian section as a section of a twisted
weight-five line.  The line, section, line comparison, character, and
transition compatibilities are therefore separate data, not consequences
of the scalar product.
\end{proof}

\subsubsection*{First-order Dirac block decomposition}

\begin{definition}[First-order Dirac block]
\label{def:first-order-dirac-block}
For $\gamma\in\Gamma_R^{\Pi,+}$, write $P^{\Pi,+}_{R,\gamma}$ for the
cusp-positive primitive piece in degree $\gamma$ and put
$x^\gamma=q^nr^ls^m$ for $\gamma=(n,l,m)$.  The two-term
Pfaffian/Koszul block on
$P^{\Pi,+}_{R,\gamma}\oplus(P^{\Pi,+}_{R,\gamma})^\vee[1]$ is the
skew operator
\[
J_\gamma=
\begin{pmatrix}0&1-x^\gamma\\-(1-x^\gamma)&0\end{pmatrix},
\qquad
\operatorname{Pf}(J_\gamma)=(1-x^\gamma)^{\sdim P^{\Pi,+}_{R,\gamma}}.
\]
For $\gamma\notin\Gamma_R^{\Pi,+}$, set $P^{\Pi,+}_{R,\gamma}=0$.
The first-order block of $\mathfrak D_X^{D_0}$ at height $R$ is the
direct sum
\[
\mathfrak D_{X,R}^{D_0}=\bigoplus_{\gamma\in\Gamma_R^{\Pi,+}\setminus\{0\}}
J_\gamma,
\]
acting on
$\bigoplus_\gamma P^{\Pi,+}_{R,\gamma}\oplus
(P^{\Pi,+}_{R,\gamma})^\vee[1]$.
\end{definition}

The principal symbol of $\mathfrak D_{X,R}^{D_0}$ with respect to the
finite Chevalley--Eilenberg/Koszul filtration is the normal-ordered
Hall bracket of \S\ref{subsec:hall-correspondences}: the Pfaffian
line is the determinant detector, and the symbol carries the Lie
data.  The phrase \emph{first-order} means the Pfaffian/primitive
object before scalar squaring; no analytic Dirac operator on a
constructed state space is asserted.

\subsubsection*{The Pfaffian product at finite stage}

\begin{proposition}[Pfaffian product at finite stage;
\textit{conditional on \((P_{\mathrm{fin}})\)}]
\label{prop:pfaffian-product-finite-stage}
Suppose the orientation lane $(\textup{O1},\textup{O1}^+)$ has
descended to the reduced quotient, the finite geometric Pfaffian datum
\((P_{\mathrm{fin}})\) of
Definition~\ref{def:finite-geometric-pfaffian-datum} is supplied, and
the support and parity residuals
\[
\mathfrak o^{\mathrm{supp}}_{R,\gamma}=
\begin{cases}
P^{\Pi,+}_{R,\gamma},&a_\Delta(\gamma)=0\\
0,&a_\Delta(\gamma)\ne0
\end{cases},
\qquad
\mathfrak o^{\mathrm{par}}_{R,\gamma}
=p^{\bar0}_{R,\gamma}-p^{\bar1}_{R,\gamma}-a_\Delta(\gamma)\in\Z,
\]
together with the leading-coefficient residual
$\mathfrak o^{\mathrm{lead}}_R=\operatorname{lc}_{\mathfrak c_\infty}
(\operatorname{pf}_{X,R})-64$, all vanish on the test window, where
$a_\Delta(\gamma)$ is the Borcherds exponent and
$p^{\bar p}_{R,\gamma}=\dim P^{\Pi,+}_{R,\gamma,\overline p}$.  Then
the Pfaffian section at finite stage has expansion
\[
\operatorname{pf}_{X,R}|_{\mathfrak c_\infty}
=64q^{1/2}r^{1/2}s^{1/2}
\prod_{\substack{\gamma\in\Gamma_R^{\Pi,+}\\\gamma\ne 0}}
(1-x^\gamma)^{\sdim P^{\Pi,+}_{R,\gamma}}.
\]
The cusp-positive Pfaffian half is recorded; negative root
representatives and the Cartan/imaginary completion enter the Lie
recognition statement of \S\ref{subsec:primitive-recognition}.
\end{proposition}

The support residual is a super-vector-space residual, not a $K_0$
class: its vanishing excludes invisible even--odd cancelling
primitive pieces in target-zero degrees.  The parity residual is a
signed superdimension residual.  These finite-stage residuals on the
test window do not by themselves prove compact source coverage; the
inclusion $A_R\subset\Gamma_R^{\Pi,+}$ of active target degrees in the
source image is itself a residual condition, not a consequence of
the Borcherds product.

\subsubsection*{Local Pfaffian wall on type-II walls}

The (O2) obstruction concerns the local Pfaffian normal form on the
three simple type-II walls $\Hpl_{\delta_i}$, $i=1,2,3$, of the
Borcherds--Cartan datum $\mathscr B_{\Delta_5}$.

\begin{definition}[Rank-one type-II Pfaffian crossing]
\label{def:rank-one-type-ii-crossing}
For a simple type-II root $\delta$, let $T_\delta$ be a real analytic
tangent chart with trivial $s_\delta$-action, set
$U_\delta=T_\delta\times(-\varepsilon,\varepsilon)$ with
$s_\delta(t,u)=(t,-u)$, and let
$\mathscr L^{\mathrm{tan}}_\delta$ be an oriented tangent Pfaffian
line with $s_\delta$-invariant unit $\upsilon_\delta$.  The rank-one
odd Pfaffian crossing is the skew family
\[
D_u=\begin{pmatrix}0&u\\-u&0\end{pmatrix},
\qquad K^{\mathrm{cross}}_\delta=
[\R\xrightarrow{u}\R],
\]
in odd Pfaffian parity.  The total local Pfaffian section is
\[
\operatorname{pf}_\delta=\upsilon_\delta\,u,
\qquad
s_\delta^*\operatorname{pf}_\delta=-\operatorname{pf}_\delta,
\qquad
s_\delta^*(\operatorname{pf}_\delta^2)=\operatorname{pf}_\delta^2.
\]
\end{definition}

\begin{proposition}[Rank-one crossing gives the type-II local sign;
\textit{proved}]
\label{prop:rank-one-crossing-o2-sign}
For $i=1,2,3$, the rank-one crossing of
Definition~\ref{def:rank-one-type-ii-crossing} forms the finite
generic type-II local-sign datum with
$N^{\mathrm{Pf}}_{\delta_i}=1$, and at finite HN height $R$, with
successor maps preserving the tangent charts, the wall coordinate
$u_{\delta_i,R}$, and the invariant units, the (O2) residual entries
\[
\mathfrak o^{\mathrm{chart}}_{\delta_i,R},\
\mathfrak o^{\mathrm{wall}}_{\delta_i,R},\
\mathfrak o^{\mathrm{split}}_{\delta_i,R},\
\mathfrak o^{\mathrm{unit}}_{\delta_i,R},\
\mathfrak o^{\mathrm{rank}}_{\delta_i,R},\
\mathfrak o^{\mathrm{tr}}_{\delta_i,R'R}
\]
vanish for these local models, and
$\varepsilon^{\mathrm{loc}}_{\delta_i,R}=-1$.
\end{proposition}

\begin{proof}
The fixed locus of $s_\delta(t,u)=(t,-u)$ is $u=0$, so the wall and
the tangent fixed locus agree.  The determinant factor is split: the
tangent line is $\mathscr L^{\mathrm{tan}}_\delta$, and the normal
factor is $K^{\mathrm{cross}}_\delta$.  The unit $\upsilon_\delta$ is
invariant by construction.  The skew matrix has Pfaffian $u$, so the
normal-mode rank is one.  Reflection sends $u\mapsto -u$ and fixes
$\upsilon_\delta$; hence the local Pfaffian section changes sign and
its square is fixed.  The transition hypothesis is exactly
preservation of these data under successor maps.
\end{proof}

\subsubsection*{Half-Hilbert orbit and the type-II sign sum}

The (O2) obstruction also requires a wall \emph{atlas}: at finite
HN height $R$, each retained component and boundary stratum of
$\Hpl_{\delta_i}$ needs a chart of the rank-one form, and the
normal Pfaffian units, tangent/normal splittings, quotient
orientations, and protected integration maps must agree on chart
overlaps.  A single generic chart computes only a local sign.

The finite O2 wall-atlas proof packet is the table system
\[
\mathrm{wall\_objects},\quad
\mathrm{wall\_charts},\quad
\mathrm{overlaps},\quad
\mathrm{orbit\_transport},\quad
\mathrm{half\_hilbert\_orbit},\quad
\mathrm{compatibilities},\quad
\mathrm{scalar\_separation}.
\]
Its rows must carry the retained wall objects, rank-one chart
coordinates, Cech/cochain overlap equations, Weyl transport, the
twelve-coordinate half-Hilbert orbit, and the component, boundary,
Thom--Sebastiani, HN-transition, quotient-orientation, and protected
integration identities.  The scalar-separation rows record
\(4096=2^{12}=64^2\) without identifying the OP scalar normalization
or the Maass character value with wall-atlas data.

The half-Hilbert orbit of the type-II reflection group on the
relevant retained stratum carries a sign sum of size $4096=2^{12}$,
counting the supplied finite-window normal Pfaffian signs across the
twelve retained components and boundary strata of the type-II wall
configuration; this sum is the signature of the (O2) wall
normal-form datum.  Local rank-one charts contribute the local
factors; the chosen Weyl lifts convert the retained global system
into the orientation character on the
$W^{(2)}(\Lambda_{II}^{2,1})$-orbit.  The full normal-form
theorem with overlap, boundary, and component coherence is relative
to the retained \((O2_{\mathrm{atlas}})\) datum and
Theorem~\ref{thm:G-O2-orbit-transport} of
Appendix~\ref{app:G-O2-discharge}; the local model is the finite
rank-one wall factor.

\begin{proposition}[Generic local sign and the orientation character]
\label{prop:generic-local-sign-orientation-character}
Suppose (O1)$^+$ supplies a normalized lift
$\tau_{i,R,S}:o_{R,S}\to s_{\delta_i}^* o_{R,S}$ of the simple
type-II reflection on every retained orbit, and the (O2)/hybrid
residual for $\delta_i$ vanishes.  In a local chart with
$K^{\mathrm{nor}}_{\delta_i,R}\simeq[\R\xrightarrow{u_i}\R]$, with
normal Pfaffian unit $\upsilon_i$ and $s_{\delta_i}(\upsilon_i)
=\upsilon_i$, let $\ell_{\delta_i,R,S}$ be the path in the reduced
quotient wall-complement carrying $u_i=\epsilon$ to $u_i=-\epsilon$,
closed by $\tau_{i,R,S}$.  Then
\[
M_{\ell_{\delta_i,R,S}}(\upsilon_i u_i)=\upsilon_i(-u_i)
=-\upsilon_i u_i,
\qquad
\epsilon_{o,R}(s_{\delta_i})=(-1)^{N^{\mathrm{Pf}}_{\delta_i,R}}.
\]
For the rank-one type-II model
$\epsilon_{o,R}(s_{\delta_i})=-1=\nu_{\Delta_5}(s_{\delta_i})$.  The
global statement
$\epsilon_o|_{W^{(2)}(\Lambda_{II}^{2,1})}=\nu_{\Delta_5}|
_{W^{(2)}(\Lambda_{II}^{2,1})}$
is the chamber-permuting half of the Pfaffian--Dirac orientation
identity, conditional on the (O2) wall atlas.
\end{proposition}

\subsubsection*{Pfaffian normalization and trace}

The supplied Pfaffian-to-automorphic comparison
$\iota_{\mathrm{aut}}$ then gives, on the imported Borcherds product
side,
\[
\iota_{\mathrm{aut}}\bigl(\operatorname{Pf}_{\mathrm{prot}}
(\mathfrak D_X)\bigr)=\mathcal D_X=\Delta_5,
\]
with cusp-side expansion
\[
\operatorname{pf}_X|_{\mathfrak c_\infty}
=64q^{1/2}r^{1/2}s^{1/2}\prod_{\gamma=(n,l,m)\in
\Gamma_{\mathrm{eff}}}(1-q^nr^ls^m)^{f(nm,l)},
\]
matching the imported Borcherds--Gritsenko--Nikulin product.  The
scalar branch obtained after taking the protected trace of the
inverse squared Pfaffian section is the Oberdieck--Pixton
normalization
\[
Z^X_{\mathrm{OP}}=-4096\,\Delta_5^{-2},
\]
where the negative sign is the OP scalar branch and
\(4096=2^{12}=64^2\) is the scalar equality between the retained
half-Hilbert sign count and the squared theta-leading coefficient.
Neither number defines the orientation character.

\subsubsection*{Comparison with the Pfaffian identity}

The Vol II primitive-to-shadow comparison singles out the protected Pfaffian
$\operatorname{Pf}_{\mathrm{prot}}(\mathfrak D_X)=\Delta_5$ as the
target identity for the retained Pfaffian datum; the \(K3\times E\)
side supplies the Pfaffian line and scalar branch data.
BBDJS \cite{BBDJS2015} provides the orientation-line and Pfaffian
calculus on the cosection-reduced d-critical heart; Joyce--Upmeier
\cite{JoyceUpmeier} provides multiplicativity under reduced
extension correspondences; Joyce--Tanaka--Upmeier
\cite{JoyceTanakaUpmeier} extends the framework to $4$-fold
analogues but is not used directly here.  The constant \(4096=2^{12}\)
is fixed at the orientation level by the retained half-Hilbert wall
atlas and at the scalar level by \(64^2\); the local sign \(-1\) is
fixed by the rank-one Pfaffian wall computation.  The
Oberdieck--Pixton minus sign is a scalar branch convention, not the
source of the orientation character.  No
identification of the Pfaffian $\Delta_5$ with a compact microscopic
state space is asserted: \(\Delta_5\) is the protected Pfaffian section,
the scalar trace is \(\Delta_5^{-2}\), and the operator-level datum is the
pro-object $\mathfrak D_X^{\mathrm{DI}}$, defined at the Pfaffian and
orientation level relative to the retained data of
Appendix~\ref{app:obstruction-discharges}.

The Pfaffian recovers $\Delta_5$ at the target scalar level, but
the cross-level recognition that identifies the primitive Hopf
algebra of \S\ref{subsec:hall-correspondences} with the Borcherds--Kac--Moody
target $\autcor$ requires a separately constructed source coalgebra
and an explicit chiral-Koszul cobar comparison.  The source coalgebra
$C_X$ is the bar of the augmented binary/two-step hybrid correspondence
algebra; it becomes the bar of a hybrid factorization algebra only
after the higher-coloured residual vanishes.  The cobar map
$\Theta_{\mathrm{Kos},R}$ is the quasi-isomorphism to the imported
current envelope $\mathsf{Den}_{\Delta_5,E}$.  This Koszul lane
$(\mathrm L 5)$ is the content of
\S\ref{subsec:koszul-source}.
